% Source: 05 - Tue Oct 6 2026.pdf, page 6. Content remains in source order.
\sourcepage{05}{6}
\section{Reductions by restriction}

These reductions are used to prove that an NPC problem stays NPC even if we restrict the set of admissible instances.
We will find the smallest subset of instances of a NPC problem for which the problem remains in NPC. 

\subsection{Restriction on Formula satisfiability}
\begin{notesdefinition}
\textbf{DEF} Given a set of boolean variables $\{x_1,x_2,\ldots,x_n\}$:
\begin{itemize}
\item A literal is either a variable $x_i$ or its negation $\neg x_j$ ($\bar x_j$). We will use $y$ to denote a literal.
\item A clause is a disjunction of 3 distinct literals, e.g. $C_i=x_1\lor\bar x_3\lor x_5$.
\item A 3-CNF formula $\phi(x_1,\ldots,x_n)$ is a conjunction of clauses built on the $n$ variables.
\end{itemize}
E.g. $\phi(x_1,x_2,x_3,x_4)=(x_1\lor\bar x_2\lor x_3)\land(\bar x_1\lor x_2\lor\bar x_3)$.

\end{notesdefinition}

\subsection{3-Conjunctive Normal Form -SAT}
We are restricting the set of legal inputs of Formula Satisfiability to 3 CNF.

\begin{notesproblem}
\text{I:} & \langle\phi(x_1,x_2,\ldots,x_n)=\bigwedge_{i=1}^m C_i\rangle\text{, }\phi\text{ 3-CNF formula.}\\
\text{Q:} & \text{Is }\phi(\vec x)\text{ satisfiable?}
\end{notesproblem}

\section{3-CNF-SAT is in NPC}

We will prove that $\mathrm{3\text{-}CNF\text{-}SAT}\in\classNPC$ : Even restricting the shape of the formula to a very simple one, the problem stays difficult.

\subsection{Step 1 : Proove that 3-CNF-SAT is in NP }
\textbf{Homework : }
$\mathrm{3\text{-}CNF\text{-}SAT}\in\classNP$ (trivial, exercise).

\subsection{Step 2 : Proove that 3-CNF-SAT is in NPH }
$\mathrm{3\text{-}CNF\text{-}SAT}\in\classNPH$:

We reduce from BC-SAT, using a similar idea to $\mathrm{BC\text{-}SAT}\polyred\mathrm{SAT}$.
